\section{A concrete next target: the conductor-23 Hecke function}
\label{sec:hecke23}

The most useful next zero-location target is a single, fully specified
degree-two function. Let
\[
 K=\Q(\sqrt{-23}),\qquad \Cl(K)\simeq C_3,
\]
and let $\theta$ be either nontrivial character of $\Cl(K)$. Complex
conjugation of ideals preserves their norms and inverts their classes, so
$L_K(s,\theta)=L_K(s,\theta^{-1})$. Denote the common function by
$L_\theta(s)$. It is the $L$-function of
\begin{equation}\label{eq:eta23}
 f(z)=\eta(z)\eta(23z)
 =q\prod_{n\geq1}(1-q^n)(1-q^{23n}),\qquad q=e^{2\pi i z}.
\end{equation}
This is a weight-one cusp form of level $23$, with quadratic character
$\chi_{-23}$. Its Artin representation $\rho$ is the standard
two-dimensional representation of $S_3$, with conductor $23$ and root
number $+1$. A cubic stem field is defined by $x^3-x^2+1$, of
discriminant $-23$. These identifications and the class group can be checked
against the LMFDB records cited here.\cite{lmfdbQuadratic,lmfdbArtin,lmfdbForm}

This target is outside the five public Dirichlet and zeta statements in the
repository. Its analytic continuation and modular interpretation already
exist; the proposed extension concerns the locations of its zeros. We do
not prove that zero-location statement, and we make no exhaustive claim
about the absence of other recent work on it. The value of fixing this
example is that every proposed transfer can be tested against a precise
Euler product, rather than against an unspecified class of automorphic
functions.

\subsection{An exact obstruction to enlarging the rational modulus}

Let $H$ be the splitting field of $x^3-x^2+1$. The character of $\rho$ has
the following values:
\[
\begin{array}{c|ccc}
 \text{Frobenius class in }S_3&1&(12)&(123)\\ \hline
 \operatorname{tr}\rho&2&0&-1.
\end{array}
\]
The current proof's rational congruence conditions cannot encode this
character merely by replacing the modulus with a larger integer.

\begin{proposition}[Finite-modulus obstruction]\label{prop:modulus-obstruction}
For every integer $q\geq1$, infinitely many unramified primes
$p\equiv1\pmod q$ have $\operatorname{tr}\rho(\operatorname{Frob}_p)=2$,
and infinitely many have
$\operatorname{tr}\rho(\operatorname{Frob}_p)=-1$.
Consequently these traces are not a function of $p$ modulo any fixed
integer, even after deleting finitely many primes. In particular they
are not a finite linear combination of Dirichlet characters.
\end{proposition}

\begin{proof}
Put $C=\Q(\zeta_q)$. The extension $H\cap C$ is abelian and Galois over
$\Q$. The largest abelian quotient of $S_3$ is its sign quotient, so
\[
 H\cap C\subseteq H^{A_3}=\Q(\sqrt{-23}).
\]
Both the identity and a $3$-cycle act trivially on this intersection.
The pairs $(1,1)$ and $((123),1)$ therefore define elements of
$\operatorname{Gal}(HC/\Q)$, viewed as the fiber product over
$\operatorname{Gal}(H\cap C/\Q)$. The Chebotarev theorem gives infinitely
many primes with each of their conjugacy classes. The cyclotomic
component is the identity, giving $p\equiv1\pmod q$, while the two traces
are $2$ and $-1$. A finite list of Dirichlet characters has a common
modulus, on whose identity residue class every character equals $1$.
No fixed linear combination can take both displayed values there.
\end{proof}

The proposition also prevents a finite product or quotient of Dirichlet
$L$-functions from having this Euler product: compare the prime
coefficients of their logarithmic derivatives. It is an obstruction to
that particular reduction, not a statement of logical independence from
Dirichlet GRH.

There is a related structural obstruction to changing only the auxiliary
field. The current field $E=\Q(\sqrt{-d},\zeta_{q_1})$ is abelian over
$\Q$, so it cannot contain $H$. Even replacing it with a general Galois
CM field would not permit that inclusion. Complex conjugation is central
in the Galois group of a Galois CM field, whereas it is a noncentral
transposition in $\operatorname{Gal}(H/\Q)$. A central element would
restrict centrally to the Galois subfield $H$, a contradiction. An
extension therefore needs relative or adelic class information, rather
than this direct field substitution.

\subsection{A positive selector and a complete analytic bridge}

For $\Re s>1$, define the logarithmic prime-ideal series
\begin{equation}\label{eq:class-selector}
 A_0(s)=
 \sum_{\substack{\mathfrak p\text{ prime},\ k\geq1\\
                         {[\mathfrak p]^k=1}}}
 \frac{\log N\mathfrak p}{(N\mathfrak p)^{ks}}.
\end{equation}
The condition is on the class of the \emph{prime power}; a nonprincipal
prime ideal contributes whenever its class is killed by $k$. Character
orthogonality on $C_3$ gives the exact identity
\begin{equation}\label{eq:class-logderiv}
 A_0(s)=-\frac13\left(
       \frac{\zeta_K'(s)}{\zeta_K(s)}
       +2\frac{L_\theta'(s)}{L_\theta(s)}\right).
\end{equation}
All coefficients in \eqref{eq:class-selector} are nonnegative. This
positivity is compatible with positive multiplicities of all zeros in
\eqref{eq:class-logderiv}; no companion function occurs with the opposite
logarithmic-derivative sign.

Write the Laurent expansion at $1$ as
\[
 A_0(s)=\frac{1}{3(s-1)}+a_0+O(s-1),\qquad c=a_0-\frac13.
\]
Define the smoothed remainder
\begin{equation}\label{eq:class-remainder}
 R_0(x)=
 \sum_{\substack{N\mathfrak p^k\leq x\\{[\mathfrak p]^k=1}}}
 \log N\mathfrak p\left(\frac1{N\mathfrak p^k}-\frac1x\right)
 -\frac13\log x-c,\qquad x\geq1.
\end{equation}
Both the subtraction constant and the smoothing normalization matter.
For $n\geq1$ and $\Re s>1$,
\[
 \int_n^\infty\left(\frac1n-\frac1x\right)x^{-s}\,dx
       =\frac{n^{-s}}{s(s-1)}.
\]
Absolute convergence allows termwise integration, giving
\begin{equation}\label{eq:class-mellin}
 M(s):=\int_1^\infty R_0(x)x^{-s}\,dx
  =\frac{A_0(s)}{s(s-1)}
       -\frac1{3(s-1)^2}-\frac{c}{s-1}.
\end{equation}
The singularity at $s=1$ is removable: the Laurent coefficient of
$(s-1)^{-1}$ in the first term is precisely $a_0-1/3=c$.

If $m_K(\rho)$ and $m_\theta(\rho)$ are the multiplicities of a zero
$\rho$ of the two functions, with a multiplicity set to zero when that
function does not vanish, then for $0<\Re\rho<1$
\begin{equation}\label{eq:class-residue}
 \operatorname*{Res}_{s=\rho} M(s)
 =-\frac{m_K(\rho)+2m_\theta(\rho)}{3\rho(\rho-1)}.
\end{equation}
Coincident zeros add their multiplicities. For a real
$\rho\in(1/2,1)$ the residue is strictly positive. Under the already
available Dirichlet GRH, $\zeta_K=\zeta L(s,\chi_{-23})$ has no zero
to the right of $1/2$, and the numerator there simplifies to
$2m_\theta(\rho)$.

\begin{theorem}[One-sided ideal-class bridge]\label{thm:hecke-bridge}
Fix $1/2<\beta<1$. Suppose that there exist $C\geq0$ and $X_0\geq1$
such that
\begin{equation}\label{eq:hecke-missing-bound}
 R_0(x)\leq Cx^{\beta-1}\qquad(x\geq X_0).
\end{equation}
Then neither $\zeta_K(s)$ nor $L_\theta(s)$ has a zero with
$\Re s>\beta$. If \eqref{eq:hecke-missing-bound} holds for every
$\beta>1/2$, both functions satisfy their Riemann hypotheses.
\end{theorem}

\begin{proof}
For $x\geq X_0$ set $F(x)=Cx^{\beta-1}-R_0(x)\geq0$. Modify $F$ on
$[1,X_0]$ to make it nonnegative everywhere. This changes its Mellin
transform by an entire function, because the correction is integrable
and compactly supported in $[1,\infty)$. Initially for $\Re s>1$, its
transform is consequently
\begin{equation}\label{eq:positive-mellin}
 H(s)=\frac{C}{s-\beta}-M(s)+J(s),
 \qquad J\text{ entire}.
\end{equation}
Absolute convergence of the series defining $A_0$ proves absolute
convergence of this integral for $\Re s>1$. Thus the abscissa of
convergence $\sigma_c$ is at most $1$, or is $-\infty$.

Suppose first that $\sigma_c>\beta$. Landau's theorem for Mellin
transforms of nonnegative functions says that $H$ has a singularity
at the real point $\sigma_c$. The standard analytic properties of
finite-order Hecke functions, together with \eqref{eq:class-mellin},
give a meromorphic continuation of \eqref{eq:positive-mellin} to
$\Re s>\beta$. The point $1$ is removable. Any real singularity in
$(\beta,1)$ must therefore come from a real zero of $\zeta_K$ or
$L_\theta$. By \eqref{eq:class-residue} its residue in $H$ is negative.
As real $s$ decreases to $\sigma_c$ from the right,
\eqref{eq:positive-mellin} would then tend to $-\infty$. But on that
interval it is the integral of the nonnegative function $F(x)x^{-s}$,
and cannot be negative. This contradiction proves
$\sigma_c\leq\beta$, including the possibility $\sigma_c=-\infty$.

The integral defining $H$ is now holomorphic on $\Re s>\beta$.
Every zero of either target function in that region would produce,
by \eqref{eq:class-residue}, a nonzero pole in
\eqref{eq:positive-mellin}. Its numerator is a positive sum of
multiplicities, so it cannot cancel. There are consequently no such
zeros. Letting $\beta$ decrease to $1/2$ gives the right-half-plane
assertion. The functional equations give the reflected assertion
on the left of the critical line.
\end{proof}

This proof is a specialization of classical Landau oscillation
technology, closely related to the one-sided philosophy used in the
repository and to Suzuki's weighted-sum criteria.\cite{suzuki,repo}
It is derived explicitly here to fix the arithmetic target and the
residue signs; no literature-priority claim is attached to it. It is
also a conditional bridge: we have \emph{not} established
\eqref{eq:hecke-missing-bound}.

\subsection{The missing arithmetic construction and a finite benchmark}

The missing estimate is specific. The existing short/enough/long
construction measures rational prime sums weighted by
$\Lambda(n)(1-\chi_d(n))$ and a rational congruence condition.
The new construction would have to make its finite-place divisibilities
measure \eqref{eq:class-selector}, while preserving the nearly spanning
family, the denominator control, and the cancellation of main heights.
Proposition~\ref{prop:modulus-obstruction} explains why the required class
selector is not supplied by changing a modulus.

There is a route closer in shape to the existing proof: introduce an
auxiliary quadratic Hecke character $\eta$ and weight each summand of
\eqref{eq:class-selector} by $1-\eta(\mathfrak p)^k$. Positivity then
survives, but the logarithmic derivative becomes a difference of
untwisted and twisted functions. One must additionally choose a twist
whose companion functions do not vanish at the proposed offending
zero. The Dirichlet-specific noncancellation input cannot be replaced
by a name change. The untwisted selector used above avoids this
analytic cancellation issue at the price of a different arithmetic
construction.

The supplied exact benchmark fixes all local normalizations. For a
rational prime $p$, the coefficient of $p^{-ks}$ in $A_0(s)$, divided
by $\log p$, is as follows:
\[
\begin{array}{c|c|c}
 \text{rational prime type}&\text{allowed }k&\text{coefficient}/\log p\\ \hline
 \operatorname{Frob}_p=1&\text{all }k&2\\
 \operatorname{Frob}_p\text{ a transposition}&2\mid k&2\\
 \operatorname{Frob}_p\text{ a }3\text{-cycle}&3\mid k&2\\
 p=23&\text{all }k&1
\end{array}
\]
All other entries are zero. The transposition line uses a prime ideal
of norm $p^2$, not $p$. At $23$ the cubic factorization is
$(x-15)(x-16)^2$ modulo $23$; its standard Artin local factor is
$(1-23^{-s})^{-1}$.

The Python benchmark independently computes the coefficients of
\eqref{eq:eta23} and the Euler product obtained from cubic splitting.
All coefficients through $n=2000$ agree, covering $303$ primes. It also
checks \eqref{eq:class-logderiv} on twelve powers of each tested prime,
for $3636$ exact coefficient checks, with the ramified prime treated
separately. These computations verify finite identities and conventions.
They give no proof concerning zeros, their locations, or their
multiplicities.

Practical milestones are therefore: first, expose the ideal-class
Euler product and this analytic bridge as independent statements;
second, construct a class-sensitive divisibility model and identify its
short and nearly spanning families; third, prove a matching lower
height bound with losses that remain small after division by the
number of lifts; finally, verify that the surviving error has the
exponent in \eqref{eq:hecke-missing-bound}. The second and third steps
are the substantial new mathematics. A computation of low zeros or
the already known functional equation does not replace them.

\section{Selmer control: a reusable theorem with essential quantifiers}
\label{sec:selmer-transfer}

Another substantial transfer is available before the final GRH theorem.
In the compact edition, the declaration
\texttt{exists\_anticyclotomic\_twist} in
\texttt{New/Selmer/TwistTheorem.lean} starts with finitely many
representations satisfying \texttt{IsBalanced}, finitely many
interpolated value families, a fixed torsion precision $\ell^M$, a
tolerance $\varepsilon>0$, and a finite power range $1\leq d\leq m$.
It constructs an auxiliary set of primes $I$ and one character $\psi$
such that all these powers simultaneously have controlled local cost,
controlled increases in the valuations of the value families, and
bounded normalized lengths of the specified Selmer modules modulo
$\ell^M$, after sufficiently enlarging the coefficient ring.
The construction uses finite-order anticyclotomic characters. It is
formulated in the repository's \texttt{UnramCMField} setting, with its
prescribed local relations; applying it to a new field or representation
also requires establishing those hypotheses.\cite{repo}

This simultaneous control is relevant to congruence denominators,
deformation problems, and arithmetic nonvanishing questions. It is not
a consequence of applying the final GRH statement. Shrinking Selmer
groups through carefully chosen local conditions has substantial
precedent, including the work of Mazur--Rubin and
Fakhruddin--Khare--Patrikis.\cite{mazurRubin,fkp} The useful extraction
here is the actual combined interface, with its hypotheses retained.

\subsection{The nonzero seed is a hypothesis}

The definition \texttt{IsFamily} in
\texttt{New/Selmer/Corepresentable.lean} explicitly contains the field
\[
 \texttt{val\_empty\_ne\_zero}:
       \operatorname{val}(\varnothing,1)\ne0.
\]
The conclusion bounds the valuation of a twisted value by the
valuation of this seed plus a finite loss. It therefore preserves
nonvanishing of an existing seed. It does not establish nonvanishing
of an initially zero central value, nor does the statement provide
the finite seed valuation needed for that application. Any use toward
central-value nonvanishing must first construct the relevant family
and prove its seed hypothesis.

\subsection{Fixed precision is not a full Selmer-group bound}

The precision $M$ is chosen before $I$ and $\psi$. The assertion has
the schematic shape
\[
 \forall M\ \exists(I_M,\psi_M)\quad
       \text{controlled Selmer length modulo }\ell^M.
\]
It does not supply one twist that works at all precisions. The following
countermodel makes that distinction independent of Selmer terminology.

\begin{proposition}[A quantifier countermodel]\label{prop:selmer-quantifiers}
For a fixed prime $\ell$, define
\[
 S_M(n)=\Z/\ell^{\max(M-n,0)}\Z,
 \qquad M,n\geq1,
\]
where exponent zero gives the zero module. Then for every $M$ some
$n$ has $\operatorname{length}_{\Z_\ell} S_M(n)=0$, but for every
fixed $n$ these lengths are unbounded as $M\to\infty$.
\end{proposition}

\begin{proof}
The length is exactly $\max(M-n,0)$. Choose $n=M$ for the first
assertion and let $M$ increase with $n$ fixed for the second.
\end{proof}

This is not a counterexample to the repository theorem. It disproves
an unjustified interchange of its quantifiers. There are further
identification requirements: the auxiliary ramification set can vary;
the local relations include additional evaluation conditions; and
the bound concerns normalized length after coefficient extensions.
These modules cannot be identified without proof with one fixed
integral Bloch--Kato Selmer group or with the full
Tate--Shafarevich group of an elliptic curve.

A practical next theorem would retain a fixed representation and
explicit local conditions, and show that the chosen twists and
coefficient extensions admit a compatible system with a uniform
bound as $M$ increases. Alternatively, one can use the existing
finite-precision statement directly for a finite congruence problem,
where that precision is the intended endpoint. The latter is a
concrete application of the available theorem. General
Bloch--Kato, finiteness of Tate--Shafarevich groups, and simultaneous
central nonvanishing remain beyond the conclusions established in
this report.
